\section{Bipedal Stance Evasion (v2)} \label{sec:bipedal}

\subsection{Formulation}
\begin{table}
  \caption{Variables and Parameters}
  \centering
  \begin{tabularx}{0.7\textwidth}{rcl}
  \toprule
    Symbol & Value & \multicolumn{1}{X}{Description} \\
  \midrule
    $d(\bm q)$ & \textit{\cref{eq:limbo:distance:ellipsoid,eq:limbo:distance:smin}} & object robot distance \\
    $d_\text{safe}$ & $\SI{0.05}{\newton\metre}$ & object robot safety distance \\
  \midrule
  PoS & \textit{\cref{eq:bipedal:constraints:stability}} & \acrlong{pos} \\
    $\gamma$ & $0.8$ & \acrshort{pos} safety Limit \\
    ZMP & \textit{\cref{eq:bipedal:constraints:stability}} & \acrlong{zmp} \\
  \midrule
    $\bm q$ & \textit{\cref{eq:bipedal:dyn:x}} & Joint Positions \\
  $\bm{\tau}$ & \textit{\cref{eq:bipedal:dyn:u}} & Joint Torques \\
    \midrule
  $\bm{\underline q},\ \bm{\overline q}$ & \textit{\cref{fig:bipedal:results:limits}} & Joint Limits \\
    $\bm{\tau}_{\text{max}}$ & \textit{\cref{fig:bipedal:results:limits}} & Max. Joint Torques \\
  \midrule
  $\bm F_L$, $\bm F_R$ & \textit{\cref{eq:bipedal:constraints:F,eq:bipedal:constraints:Fz}} & Ground Reaction Forces \\
    % $F_\text{safe}$ & $\SI{1}{\newton}$ & Friction Force Safety Margin \\
    $\mu$ & $0.5$ & Friction Coefficient \\
  \midrule
    $v_{\text{stat.}}$ & \SI{0.1}{\frac{\metre}{\second}} & Allowed Stationary Velocity \\
    $q_\epsilon$ & \SI{0.1}{\radian} & Terminal Pose Offset \\
  \midrule
    $T_f$ & $\SI{2.5}{\second}$ & Terminal Time \\
    $N$ & $120$ & Shooting Nodes \\
    $l_q$ & \textit{\cref{eq:bipedal:obj:joint}} & Lagrange Obj. Pose Term \\
    $\omega_q$ & $1$ & Lagrange Obj. Weight Pose Term \\
    $l_\tau$ & \textit{\cref{eq:bipedal:obj:joint}} & Lagrange Obj. Regularization Term \\
    $\omega_\tau$ & $10^{-4}$ & Lagrange Obj. Weight Pose Term \\
  \bottomrule
  \end{tabularx} \label{tab:bipedal:variables}
\end{table}

\begin{subequations}
  \label{eq:bipedal:dyn}
  \begin{align}
    \bm x(t) &=
    \begin{bmatrix}
      \bm{q} \\ \bm{\dot q}
    \end{bmatrix} \label{eq:bipedal:dyn:x} \\
    \bm u(t) &=
    \bm{\tau} \label{eq:bipedal:dyn:u}\\
    \bm{f}_\text{expl}(t) &=
    \begin{bmatrix}
      \bm{\dot q} \\
      \text{cons. FD}(\bm q,\bm{\dot q},\bm{\tau})
    \end{bmatrix} \label{eq:bipedal:dyn:f}
  \end{align}
\end{subequations}

\begin{subequations}
  \label{eq:bipedal:constraints}
  \begin{align}
    \min_{\bm x(\cdot),\bm u(\cdot)} \quad & \int_0^{T_f} \left(\omega_q l_\text{q}(\bm x) + \omega_\tau l_\tau(\bm{u}) \right) dt \label{bipedal:constraints:obj} \\
    \textrm{s.t.} \quad & \qquad \textbf{Dodge Constraint} \nonumber \\
    & d_\text{safe} \leq d(\bm p) & & \forall t \in [0,T_f) \label{eq:bipedal:constraints:dodge} \\
    & \qquad \textbf{Stability Constraint} \nonumber \\
    & \text{ZMP} \in \gamma \text{PoS} & & \forall t \in [0,T_f) \label{eq:bipedal:constraints:stability} \\
    & \qquad \textbf{Joint Constraints} \nonumber \\
    & \bm{\underline q} \leq \bm{q} \leq \bm{\overline q} & & \forall t \in [0,T_f) \label{eq:bipedal:constraints:joint}\\
    & -\bm{\tau}_\text{max} \leq \bm u \leq \bm{\tau}_\text{max} & & \forall t \in [0,T_f)\label{eq:bipedal:constraints:torque}\\
    & \qquad \textbf{Friction Constraints} \nonumber \\
    & 0 \leq (\mu F_j^z)^2-\|\bm{F}_j\|^2_2 & \forall j \in \{ \text{left}, \text{right} \} \qquad & \forall t \in [0,T_f) \label{eq:bipedal:constraints:F} \\
    & 0 \leq F_L^z & \forall j \in \{ \text{left}, \text{right} \} \qquad & \forall t \in [0,T_f) \label{eq:bipedal:constraints:Fz} \\
    & \qquad \textbf{Terminal Pose Constraint} \nonumber \\
    & \begin{bmatrix} \bm q_0-q_\epsilon \\ -\bm v_\text{stat.} \end{bmatrix}
    \leq \bm x(T_f) \leq
    \begin{bmatrix} \bm q_0+q_\epsilon \\ \bm v_\text{stat.} \end{bmatrix} \label{eq:bipedal:constraints:terminal}
  \end{align}
\end{subequations}

\begin{subequations}
  \label{eq:bipedal:obj}
  \begin{align}
    l_q(\bm x) &= \frac{\| \bm q-\bm q_0 \|^2_2}{(n_q \cdot \bm{q}_\text{exp}^2)^2}  \label{eq:bipedal:obj:joint} \\
    l_\tau(\bm u) &= \frac{\|\bm{\tau}\|^2_2}{(n_u \cdot \bm{u}_\text{exp}^2)^2} \label{eq:bipedal:obj:torque}
  \end{align}
\end{subequations}

\subsection{Results \& Discussion}
\begin{figure}
  \centering
  \begin{subfigure}[b]{0.45\textwidth}
    \includegraphics[width=\textwidth]{figures/06_bipedalEvasion/conv/top.pdf}
    \caption{Top View (XY)}
    \label{fig:bipedal:results:convTop}
  \end{subfigure}
  \hfill
  \begin{subfigure}[b]{0.45\textwidth}
    \includegraphics[width=\textwidth]{figures/06_bipedalEvasion/conv/iso.pdf}
    \caption{Isometric View}
    \label{fig:bipedal:results:convIso}
  \end{subfigure}
  \vfill
  % \hfill
  \begin{subfigure}[b]{0.45\textwidth}
    \includegraphics[width=\textwidth]{figures/06_bipedalEvasion/conv/front.pdf}
    \caption{Front View (YZ)}
    \label{fig:bipedal:results:convFront}
  \end{subfigure}
  \hfill
  \begin{subfigure}[b]{0.45\textwidth}
    \includegraphics[width=\textwidth]{figures/06_bipedalEvasion/conv/side.pdf}
    \caption{Side View (XZ)}
    \label{fig:bipedal:results:convSide}
  \end{subfigure}
  \caption{orange: no convergence, gray lines are not converged runs}
  \label{fig:bipedal:results:conv}
\end{figure}

\begin{figure}
  \centering
  \includegraphics[width=\textwidth]{figures/06_bipedalEvasion/plots/torqueJoint_original.pdf}
  \caption{
    Joint angles $\bm{q}$ and torques $\bm{\tau}$ for all successful dodging motions.
    We can clearly see that some of those trajectories require the robot to operate at its limits.
    Torques in general show a high usage of the available strength, indicating that the success rate is indeed limited by torque limits.
  }
  \label{fig:bipedal:results:limits}
\end{figure}

\begin{figure}
  \centering
  \includegraphics[width=0.7\textwidth]{figures/06_bipedalEvasion/plots/distance.pdf}
  \caption{Distance of traveling object throughout the evasion motions over all successfully evaded object trajectories, with the minimal distance being clearly above the safety distance to the collision hull}
  \label{fig:bipedal:results:distance}
\end{figure}

\subsubsection{Motion Classes}

\begin{figure}
  \centering
  \begin{subfigure}[b]{0.23\textwidth}
    \includegraphics[width=\textwidth]{figures/06_bipedalEvasion/motions/crouch.png}
    \caption{crouching evasion}
    \label{fig:bipedal:results:motion:crouch}
  \end{subfigure}
  \hfill
  \begin{subfigure}[b]{0.23\textwidth}
    \includegraphics[width=\textwidth]{figures/06_bipedalEvasion/motions/lean.png}
    \caption{leaning evasion}
    \label{fig:bipedal:results:motion:lean}
  \end{subfigure}
  % \vfill
  \hfill
  \begin{subfigure}[b]{0.23\textwidth}
    \includegraphics[width=\textwidth]{figures/06_bipedalEvasion/motions/inverseNeo.png}
    \caption{reverse neo evasion}
    \label{fig:bipedal:results:motion:inverseNeo}
  \end{subfigure}
  \hfill
  \begin{subfigure}[b]{0.23\textwidth}
    \includegraphics[width=\textwidth]{figures/06_bipedalEvasion/motions/neo.png}
    \caption{classic neo evastion}
    \label{fig:bipedal:results:motion:neo}
  \end{subfigure}
  \caption{
    We found, that our optimal object avoidance generator explored mainly four motion classes, for different kinds of approaching objects.
    The leaning evasion in (b) was found to be the most common one, while the reverse neo motion was leveraged in a small amount of object avoidance motions.
  }
  \label{fig:bipedal:results:motion}
\end{figure}


